\subsection{\textbf{RQ2 : What are the main challenges in type inference and
  type analysis for Python programs?}}

\begin{table*}[ht]
\centering
\small
\setlength{\tabcolsep}{4pt} % 열 간격 조정
\caption{RQ2: Table showing whether existing Python type analysis papers and
  tools support key Python features, including structural typing,
  flow-sensitive types, dynamic features, and compound data types.
  (\cmark~ is support, \xmark~ is not.)}
  \begin{tabularx}{\textwidth}{>{\centering\arraybackslash}X
                               >{\centering\arraybackslash}X
                               >{\centering\arraybackslash}X
                               >{\centering\arraybackslash}X
                               >{\centering\arraybackslash}X}
  \toprule
  \rowcolor{yellow!30}
  \textbf{Paper/Tool} & \textbf{Structural Typing} & \textbf{Flow-sensitivity} & \textbf{Dynamic Features} & \textbf{Compound DataTypes} \\
  \midrule
    \rowcolor{green!10} \cite{static2014} & \xmark & \xmark & \xmark & \cmark \\
    \rowcolor{white} \cite{static2015} & \xmark & \xmark & \xmark & \cmark \\
    \rowcolor{green!10} \cite{MaxSMT2018} & \xmark & \cmark & \xmark & \cmark \\
    \rowcolor{white} \cite{prob2016} & \xmark & \xmark & \cmark & \cmark \\
    \rowcolor{green!10} \cite{DL2018} & \xmark & \xmark & \cmark & \cmark \\
    \rowcolor{white} \cite{typewriter2020} & \xmark & \xmark & \cmark & \cmark \\
    \rowcolor{green!10} \cite{lambdanet2020} & \xmark & \xmark & \cmark & \cmark \\
    \rowcolor{white} \cite{pyter2022} & \xmark & \cmark & \cmark & \cmark \\
    \rowcolor{green!10} \cite{static2021} & \xmark & \cmark & \cmark & \cmark \\
    \rowcolor{white} \cite{hityper2022} & \xmark & \cmark & \cmark & \cmark \\
    \rowcolor{green!10} \cite{type4py2022} & \xmark & \xmark & \cmark & \cmark \\
    \rowcolor{white} \cite{pyinfer2021} & \xmark & \xmark & \cmark & \cmark \\
    \rowcolor{green!10} \cite{static2023} & \xmark & \xmark & \cmark & \cmark \\
    \rowcolor{white} \cite{pyty2024} & \xmark & \cmark & \cmark & \cmark \\
    \rowcolor{green!10} \cite{dlinfer2023} & \xmark & \cmark & \cmark & \cmark \\
    \rowcolor{white} \cite{tipical2023} & \xmark & \cmark & \xmark & \cmark \\
    \rowcolor{green!10} \cite{typeT52023} & \xmark & \xmark & \cmark & \cmark \\
    \rowcolor{white} \cite{pyinder2024} & \cmark & \cmark & \cmark & \cmark \\
    \rowcolor{green!10} \cite{quac2024} & \cmark & \cmark & \cmark & \cmark \\
    \rowcolor{white} \cite{monkeytype} & \xmark & \xmark & \cmark & \cmark \\
    \rowcolor{green!10} \cite{crosshair} & \xmark & \xmark & \cmark & \cmark \\
    \rowcolor{white} \cite{typeguard} & \xmark & \xmark & \cmark & \cmark \\
    \rowcolor{green!10} \cite{mopsa2020} & \cmark & \cmark & \cmark & \cmark \\
    \rowcolor{white} \cite{mypy} & \xmark & \xmark & \xmark & \cmark \\
    \rowcolor{green!10} \cite{pyre} & \xmark & \cmark & \cmark & \cmark \\
    \rowcolor{white} \cite{pytype} & \xmark & \cmark & \cmark & \cmark \\
    \rowcolor{green!10} \cite{pyright} & \xmark & \cmark & \cmark & \cmark \\
  \midrule
  \rowcolor{pink!30} 27 & 3 & 13 & 22 & 27 \\
  \bottomrule
  \end{tabularx}
  \label{tab:RQ2}
\end{table*}


Type analysis in Python is fundamentally more challenging than in traditional
statically typed languages due to its structural and dynamic characteristics.
In particular, (1) structural typing–based object compatibility, (2)
flow-sensitive types that change according to control flow, (3) advanced
dynamic features such as reflection and dynamic code generation, and (4) the
presence of heterogeneous containers are frequently cited as the primary
obstacles to precise type analysis.
\autoref{tab:RQ2} shows whether existing research and tools support these four key
Python features.
Each column corresponds to one of the four challenges, and the entries indicate
whether a particular paper or tool provides support.
This provides a clear overview of the strengths and limitations of current
Python type analysis approaches.

First, Python follows a duck-typing model that is closer to structural typing
than nominal typing. The meaning of an object is determined by the methods and
attributes it provides, making it difficult to accurately model objects using
only static type names. Most existing studies do not use structural typing;
instead, they perform analysis based solely on type names (nominal typing). For
this reason, current approaches cannot dynamically model Python’s types,
leading to high uncertainty in static type inference.

Second, the type of a Python variable may vary depending on the program’s
control flow. A single variable may assume different types across different
branches, making it difficult for traditional static analysis to infer a
consistent type. To address this, recent studies perform flow-sensitive
analysis, but only about 50\% of the works apply it; the rest rely on methods
derived from static-typed languages that do not account for control flow.
Flow-insensitive analysis sometimes uses techniques such as Static Single
Assignment (SSA) to represent control flow. However, increasing precision using
these methods rapidly raises computational costs. Therefore, flow-aware
analysis is essential, and balancing accuracy with performance remains a
central challenge in Python type analysis.

Third, Python’s powerful dynamic features—reflection and dynamic code
generation—fundamentally threaten the soundness of any static type analysis.
Dynamic attribute access via getattr and \_\_getattribute\_\_, runtime code
generation using eval() and exec(), dynamic module loading via importlib, and
metaclass-based class generation all defer program structure determination
until runtime. Accurately analyzing these features would require nearly full
program execution, making precise static type inference difficult within
traditional frameworks. However, as shown in \autoref{tab:RQ2}, about 80\% of studies
support some dynamic features; this is a generous count that includes partial
support. Full support for all dynamic features is not achievable in static
analysis except through execution-based approaches.

Finally, Python’s lists, sets, and dictionaries inherently allow heterogeneous
elements. A single list may contain integers, strings, and user-defined
objects, and element types may evolve during execution. Simple uniform type
models are insufficient, and more sophisticated representations—such as union
types, intersection types, or lattice-based type spaces—are required. All
surveyed studies and tools account for this feature, highlighting that
heterogeneous containers are a critically important aspect of Python type
analysis.

\begin{tcolorbox}[colback=white,colframe=blue!60,title=Answer to RQ2,fonttitle=\bfseries]
  Python type analysis exhibits high complexity due to
  (1) structural typing–based non-uniformity,
  (2) types that change continually depending on control flow,
  (3) the possibility of runtime code generation, and
  (4) heterogeneous containers.
  These characteristics make it difficult for traditional static analysis
  frameworks to achieve accurate and sound modeling, highlighting the need for
  new analysis paradigms tailored to Python’s inherently dynamic nature.
\end{tcolorbox}
